% GENERATED FILE. Edit canonical lessons, then run npm run generate:books.
% canonical-source-hash: fnv1a64:fd2bf340fffa1544
% canonical-lessons: ML-C512-karakaviyuka, ML-C512-vellakkettu, ML-C512-ponnuka, ML-C512-munnippokuka, ML-C512-valitiriccuvituka

\chapter{When the River Rises}
\label{ch:ml-when-the-river-rises}
\hlchaptermodality{512}

\begin{chapteropening}
\textbf{By the end of this chapter:} \emph{I can say to overflow its banks, waterlogging, to float up, to go under and to divert.}

Complete the last lesson of chapter 512: I can say to overflow its banks, waterlogging, to float up, to go under and to divert.
\end{chapteropening}

\section[\texorpdfstring{\ml{കരകവിയുക} (karakaviyuka)}{കരകവിയുക (karakaviyuka)}]{\ml{കരകവിയുക} (karakaviyuka) — to overflow its banks}

\label{lesson:ML-C512-karakaviyuka}



\begin{quote}
Before the new one: say the Malayalam for a cyclone, then the Malayalam for to snap and fall.
\end{quote}

\subsection*{You'll want to know: \ml{കരകവിയുക}}
\textbf{\ml{കരകവിയുക}} — \emph{karakaviyuka} — \textquotedblleft{}to overflow its banks\textquotedblright{}.

\begin{quote}
\ml{കനത്ത} \ml{മഴയിൽ} \ml{പുഴ} \ml{കരകവിഞ്ഞു}.
\end{quote}

\begin{quote}
\emph{kanatta maḻayil puḻa karakaviññu.}
\end{quote}

\begin{quote}
In the heavy rain the river overflowed its banks.
\end{quote}

Past: \ml{കരകവിഞ്ഞു}, \emph{karakaviññu}.

\ml{കര}, the bank, plus \ml{കവിയുക}, to overflow.

\begin{cousinweb}
\textbf{\ml{കരകവിയുക}} (\emph{karakaviyuka}) is a native Dravidian compound of \ml{കര}, a bank or shore (DEDR 1293), and \ml{കവിയുക}, to overflow (DEDR 1221).
\end{cousinweb}

\begin{grammarlens}[title={What you've built}]
One more word for getting around.
\end{grammarlens}

\subsection*{Your turn}
\begin{itemize}
\raggedright
  \item \emph{Say it:} \emph{karakaviyuka}
  \item \emph{Say it:} \emph{karakaviyuka}, once more
  \item \emph{Say it:} \emph{kanatta maḻayil puḻa karakaviññu.} — In the heavy rain the river overflowed its banks.
  \item \emph{Recall:} say the Malayalam for a cyclone, then the Malayalam for to snap and fall, then say \emph{karakaviyuka} again
\end{itemize}

\begin{tcolorbox}[breakable,colback=teal!4,colframe=teal!35!black,title={Before you move on}]
What does \ml{കരകവിയുക} mean? (\textquotedblleft{}To overflow its banks\textquotedblright{}.) Say it once more.
\end{tcolorbox}
\section[\texorpdfstring{\ml{വെള്ളക്കെട്ട്} (veḷḷakkeṭṭŭ)}{വെള്ളക്കെട്ട് (veḷḷakkeṭṭŭ)}]{\ml{വെള്ളക്കെട്ട്} (veḷḷakkeṭṭŭ) — waterlogging}

\label{lesson:ML-C512-vellakkettu}



\begin{quote}
Before the new one: say the Malayalam for to snap and fall, then the Malayalam for to overflow its banks.
\end{quote}

\subsection*{You'll want to know: \ml{വെള്ളക്കെട്ട്}}
\textbf{\ml{വെള്ളക്കെട്ട്}} — \emph{veḷḷakkeṭṭŭ} — \textquotedblleft{}waterlogging\textquotedblright{}.

\begin{quote}
\ml{മഴ} \ml{പെയ്താൽ} \ml{ആ} \ml{വഴിയിൽ} \ml{വെള്ളക്കെട്ട്} \ml{ഉണ്ടാകും}.
\end{quote}

\begin{quote}
\emph{maḻa peytāl ā vaḻiyil veḷḷakkeṭṭŭ uṇṭākuṁ.}
\end{quote}

\begin{quote}
When it rains, that road gets waterlogged.
\end{quote}

\ml{വെള്ളം}, water, plus \ml{കെട്ട്}, a tying up: water that cannot drain away.

\begin{cousinweb}
\textbf{\ml{വെള്ളക്കെട്ട്}} (\emph{veḷḷakkeṭṭŭ}) is a native Dravidian compound of \ml{വെള്ളം}, water, and \ml{കെട്ട്}, a tying or damming (DEDR 1147).
\end{cousinweb}

\begin{grammarlens}[title={What you've built}]
One more word for getting around.
\end{grammarlens}

\subsection*{Your turn}
\begin{itemize}
\raggedright
  \item \emph{Say it:} \emph{veḷḷakkeṭṭŭ}
  \item \emph{Say it:} \emph{veḷḷakkeṭṭŭ}, once more
  \item \emph{Say it:} \emph{maḻa peytāl ā vaḻiyil veḷḷakkeṭṭŭ uṇṭākuṁ.} — When it rains, that road gets waterlogged.
  \item \emph{Recall:} say the Malayalam for to snap and fall, then the Malayalam for to overflow its banks, then say \emph{veḷḷakkeṭṭŭ} again
\end{itemize}

\begin{tcolorbox}[breakable,colback=teal!4,colframe=teal!35!black,title={Before you move on}]
What does \ml{വെള്ളക്കെട്ട്} mean? (\textquotedblleft{}Waterlogging\textquotedblright{}.) Say it once more.
\end{tcolorbox}
\section[\texorpdfstring{\ml{പൊങ്ങുക} (poṅṅuka)}{പൊങ്ങുക (poṅṅuka)}]{\ml{പൊങ്ങുക} (poṅṅuka) — to float up; to rise (of water)}

\label{lesson:ML-C512-ponnuka}



\begin{quote}
Before the new one: say the Malayalam for to overflow its banks, then the Malayalam for waterlogging.
\end{quote}

\subsection*{You'll want to know: \ml{പൊങ്ങുക}}
\textbf{\ml{പൊങ്ങുക}} — \emph{poṅṅuka} — \textquotedblleft{}to float up; to rise (of water)\textquotedblright{}.

\begin{quote}
\ml{വെള്ളം} \ml{പൊങ്ങിയപ്പോൾ} \ml{ഞങ്ങൾ} \ml{മുകളിലത്തെ} \ml{നിലയിലേക്ക്} \ml{മാറി}.
\end{quote}

\begin{quote}
\emph{veḷḷaṁ poṅṅiyappōḷ ñaṅṅaḷ mukaḷilatte nilayilēkkŭ māṟi.}
\end{quote}

\begin{quote}
When the water rose, we moved to the upper floor.
\end{quote}

Past: \ml{പൊങ്ങി}, \emph{poṅṅi}.

For water and for things that float up in it.

\begin{cousinweb}
\textbf{\ml{പൊങ്ങുക}} (\emph{poṅṅuka}) is native Dravidian; DEDR gives the Malayalam verb as to rise up out of the water (DEDR 4471).
\end{cousinweb}

\begin{grammarlens}[title={What you've built}]
One more word for getting around.
\end{grammarlens}

\subsection*{Your turn}
\begin{itemize}
\raggedright
  \item \emph{Say it:} \emph{poṅṅuka}
  \item \emph{Say it:} \emph{poṅṅuka}, once more
  \item \emph{Say it:} \emph{veḷḷaṁ poṅṅiyappōḷ ñaṅṅaḷ mukaḷilatte nilayilēkkŭ māṟi.} — When the water rose, we moved to the upper floor.
  \item \emph{Recall:} say the Malayalam for to overflow its banks, then the Malayalam for waterlogging, then say \emph{poṅṅuka} again
\end{itemize}

\begin{tcolorbox}[breakable,colback=teal!4,colframe=teal!35!black,title={Before you move on}]
What does \ml{പൊങ്ങുക} mean? (\textquotedblleft{}To float up; to rise (of water)\textquotedblright{}.) Say it once more.
\end{tcolorbox}
\section[\texorpdfstring{\ml{മുങ്ങിപ്പോകുക} (muṅṅippōkuka)}{മുങ്ങിപ്പോകുക (muṅṅippōkuka)}]{\ml{മുങ്ങിപ്പോകുക} (muṅṅippōkuka) — to go under; to be submerged}

\label{lesson:ML-C512-munnippokuka}



\begin{quote}
Before the new one: say the Malayalam for waterlogging, then the Malayalam for to float up.
\end{quote}

\subsection*{You'll want to know: \ml{മുങ്ങിപ്പോകുക}}
\textbf{\ml{മുങ്ങിപ്പോകുക}} — \emph{muṅṅippōkuka} — \textquotedblleft{}to go under; to be submerged\textquotedblright{}.

\begin{quote}
\ml{വെള്ളപ്പൊക്കത്തിൽ} \ml{ഗ്രാമത്തിലെ} \ml{വയലുകൾ} \ml{മുങ്ങിപ്പോയി}.
\end{quote}

\begin{quote}
\emph{veḷḷappokkattil grāmattile vayalukaḷ muṅṅippōyi.}
\end{quote}

\begin{quote}
In the flood the village fields went under water.
\end{quote}

Past: \ml{മുങ്ങിപ്പോയി}, \emph{muṅṅippōyi}.

\ml{മുങ്ങി}, sinking, plus \ml{പോകുക}, to go: -\ml{പ്പോകുക} adds the sense that it happened beyond help.

\begin{cousinweb}
\textbf{\ml{മുങ്ങിപ്പോകുക}} (\emph{muṅṅippōkuka}) is a native Dravidian compound of \ml{മുങ്ങുക}, to sink under water (DEDR 4993), and \ml{പോകുക}, to go.
\end{cousinweb}

\begin{grammarlens}[title={What you've built}]
One more word for getting around.
\end{grammarlens}

\subsection*{Your turn}
\begin{itemize}
\raggedright
  \item \emph{Say it:} \emph{muṅṅippōkuka}
  \item \emph{Say it:} \emph{muṅṅippōkuka}, once more
  \item \emph{Say it:} \emph{veḷḷappokkattil grāmattile vayalukaḷ muṅṅippōyi.} — In the flood the village fields went under water.
  \item \emph{Recall:} say the Malayalam for waterlogging, then the Malayalam for to float up, then say \emph{muṅṅippōkuka} again
\end{itemize}

\begin{tcolorbox}[breakable,colback=teal!4,colframe=teal!35!black,title={Before you move on}]
What does \ml{മുങ്ങിപ്പോകുക} mean? (\textquotedblleft{}To go under; to be submerged\textquotedblright{}.) Say it once more.
\end{tcolorbox}
\section[\texorpdfstring{\ml{വഴിതിരിച്ചുവിടുക} (vaḻitiriccuviṭuka)}{വഴിതിരിച്ചുവിടുക (vaḻitiriccuviṭuka)}]{\ml{വഴിതിരിച്ചുവിടുക} (vaḻitiriccuviṭuka) — to divert (traffic)}

\label{lesson:ML-C512-valitiriccuvituka}



\begin{quote}
Before the new one: say the Malayalam for to float up, then the Malayalam for to go under.
\end{quote}

\subsection*{You'll want to know: \ml{വഴിതിരിച്ചുവിടുക}}
\textbf{\ml{വഴിതിരിച്ചുവിടുക}} — \emph{vaḻitiriccuviṭuka} — \textquotedblleft{}to divert (traffic)\textquotedblright{}.

\begin{quote}
\ml{വെള്ളക്കെട്ട്} \ml{കാരണം} \ml{പോലീസ്} \ml{ബസുകൾ} \ml{വഴിതിരിച്ചുവിട്ടു}.
\end{quote}

\begin{quote}
\emph{veḷḷakkeṭṭŭ kāraṇaṁ pōlīsŭ basukaḷ vaḻitiriccuviṭṭu.}
\end{quote}

\begin{quote}
Because of the waterlogging, the police diverted the buses.
\end{quote}

Past: \ml{വഴിതിരിച്ചുവിട്ടു}, \emph{vaḻitiriccuviṭṭu}.

\ml{വഴി}, the way, plus \ml{തിരിച്ചുവിടുക}, to turn and send.

\begin{cousinweb}
\textbf{\ml{വഴിതിരിച്ചുവിടുക}} (\emph{vaḻitiriccuviṭuka}) is a native Dravidian compound of \ml{വഴി}, a way (DEDR 5297), \ml{തിരിക്കുക}, to turn (DEDR 3246), and \ml{വിടുക}, to let go.
\end{cousinweb}

\begin{grammarlens}[title={What you've built}]
One more word for getting around.
\end{grammarlens}

\subsection*{Your turn}
\begin{itemize}
\raggedright
  \item \emph{Say it:} \emph{vaḻitiriccuviṭuka}
  \item \emph{Say it:} \emph{vaḻitiriccuviṭuka}, once more
  \item \emph{Say it:} \emph{veḷḷakkeṭṭŭ kāraṇaṁ pōlīsŭ basukaḷ vaḻitiriccuviṭṭu.} — Because of the waterlogging, the police diverted the buses.
  \item \emph{Recall:} say the Malayalam for to float up, then the Malayalam for to go under, then say \emph{vaḻitiriccuviṭuka} again
\end{itemize}

\begin{tcolorbox}[breakable,colback=teal!4,colframe=teal!35!black,title={Before you move on}]
What does \ml{വഴിതിരിച്ചുവിടുക} mean? (\textquotedblleft{}To divert (traffic)\textquotedblright{}.) Say it once more.
\end{tcolorbox}
